% s-028-p028-m002; original PDF 28
\Aphorism{s-028-p028-m002}{18}
Der Skepsis eines vornehmen Römers „was ist Wahrheit?“ wäre noch die Skepsis eines gläubigen Juden an die Seite zu setzen, etwa in die Frage gekleidet „was ist Satzung?“ Dieser Jesus, der aus der Wahrheit war und die Satzung zu erfüllen kam, musste als ein Wunder intellectueller und ethischer Naivetät angestaunt werden, von den Zöglingen einer leidlich orientirten Philosophie einerseits, die schon von der Bedingtheit alles Erkennens zu reden anfing, von den Hütern einer sehr entwickelten Moralcasuistik andrerseits, die den Menschen mit einem geradezu unentwirrbaren Dickicht von Gesetzen umstellt hatte. Jesus wies einen Ausweg aus allen Conflicten auf einmal, er improvisirte eine Lösung von verdächtigster Einfachheit, er gab ein Universalheilmittel — folglich rieth man auf Dilettantismus. Dieser Galiläer, kein Zweifel, machte es sich zu \Sperr{leicht}; er hatte die Transcendenz des Genies, das über Schwierigkeiten hinwegkommt, weil es sie nicht sieht.
\par

